%%%%%%%%%%%%%%%%%%%%%%%%%%%%%%%%%
%
%       EXERCÍCIO
%
%%%%%%%%%%%%%%%%%%%%%%%%%%%%%%%%%

\definevalorquestao

\begin{exercicioBanco}[\valorquestao]
    Admitindo que a pressão sanguínea arterial em homens siga o modelo Normal, 7 pacientes foram sorteados e tiveram sua pressão medida com os seguintes resultados: 84, 81, 77, 85, 69, 80 e 79. Teste se a média é 82 contra a alternativa de ser 78. Use $\alpha = 2\%$, onde \(\alpha\) é a probabilidade do erro do tipo I.
\end{exercicioBanco}

